\section*{Overview of the proof}
\addcontentsline{toc}{section}{Overview of the proof}
\label{sec:overview}

To estimate a prime-power sum of length \(x\), the proof composes
\(m\asymp\sqrt{x}\) theta lifts. The prime-sum remainder then appears with
a factor of \(m\) in a bound for a Gram determinant. Upper and lower bounds
for the same determinant cancel the common main term and leave a one-sided
bound for the remainder.

\subsection*{Why iterate the lift?}

In Stepanov's method for curves over finite fields, an auxiliary function
vanishes to high order at many points, and the sum of the orders of
vanishing is compared with its degree. In his announcement, Alpöge
describes iterated theta lifts over number fields as performing the role
of this multiplicity. In his account of the original proof, Zhang
also emphasizes that the leading finite-place valuation gains grow with
the number \(m\) of stages. \cite[pp.~1, 4]{alpoege} \cite[Section~4]{zhang}

In the readable edition, one estimates divisibility at finite places
after differentiating a theta kernel. For example, if \(p\) is inert
in \(K=\mathbb Q(\sqrt{-d})\), then \(1-\chi_K(p^j)\) is \(2\) for odd
\(j\) and \(0\) for even \(j\). The contribution of \(p^j\) to the local
valuation includes the cutoff weight. Summing the local gains with weight
\(2\log p\) produces the required weighted prime-power sums. A seesaw
identity connects the chain whose divisibility can be calculated to the
chain whose Gram determinant is estimated.
See Section~\ref{sec:upper-bound}.

\subsection*{The main term cancels in the determinant comparison}

Let \(c\) bound the logarithm of the squared covolume of \(n\) lifted
forms, and let \(\delta\) bound the logarithmic index required to achieve
integrality at the finite places. The index occurs as \(2\delta\)
because one change-of-coordinates matrix enters for each variable of
the Hermitian pairing. Let \(h\) be the target dimension, \(H_0\) the
common main term, and \(\mathcal R\) the total prime-sum remainder.
Write \(f\) for the number of infinite places and \(b\) for the weight
at the selected places, and set \(a=xb\) and \(E_j=fa(\log b)^j\).
The bounds take the form
\[
 h(H_0-3E_4)
 \le c+2\delta
 \le n(H_0-m\mathcal R+E_4)
     +20h\log(1+2b^2)E_3 .
\]

For the upper bound, divisibility enters the normalization of the
Poincaré map, after which one estimates the squared norms at infinity.
A second moment for translates of the map yields
\(n\ge(1-4/b^2)h\) linearly independent images of short forms.
For the lower bound, one selects a twist at each finite place and
compares the untwisted Gram determinant with one whose integrality
is controlled at that place. The height of the determinant ratio uses
Hecke \(L\)-values and bounds for denominators of projections to the
spaces of lifts. The product formula then gives the same \(H_0\).
See Sections~\ref{sec:lattices}--\ref{sec:determinant-estimates}.

Using \(n\le h\) and \(H_0\ge0\) cancels \(H_0\), and \(n\ge h/2\)
gives \(m\mathcal R\le fa(\log b)^5\). An identity using two weights
at the infinite places extracts the remainder \(R_{q_0}(x)\) in
the residue class \(1\) modulo \(q_0\). The fixed-modulus prime number
theorem controls the terms at the larger cutoff \(x_1\). The leading
term of the remaining bound is \(\varphi(q_0)(\log b)^6/m\).
Since \(m\ge\sqrt{x}\) and \(\log b\) can be chosen below an
appropriately small positive power of \(x\), this gives the one-sided
bound \(R_{q_0}(x)\le x^{-1/2+\varepsilon}\).
See Sections~\ref{sec:data} and \ref{sec:conclusion}.

\begin{editorialnote}
The working guide to the zeta case published by Rubinstein explains
the square-root scale in the original, longer proof. After
normalization by the coefficient of the prime sum, errors on scales
\(x/m\) and \(m\) balance when \(m\asymp\sqrt{x}\). The comparison and
the bound above are those of the readable edition.
\cite[Section~I.8]{rubinstein}
\end{editorialnote}

\subsection*{Why a one-sided bound excludes zeros}

The remainder \(R_{q_0}(x)\) comes from summing over
\(n\equiv1\pmod{q_0}\) with weight
\(\Lambda(n)\bigl(1-(n/d)\bigr)(1/n-1/x)\) and subtracting
\(\log x/\varphi(q_0)+c_{q_0}\). Zeros of the original
\(L\)-functions and zeros of their quadratic twists enter the
transform of this remainder with opposite residue signs.
The choice of \(d\), made with the height of the candidate zero
in mind, excludes twisted zeros in the relevant region. The
remaining residues on the real axis are nonnegative. The
one-sided bound supplies a nonnegative density for Landau's lemma,
which rules out the pole associated to the candidate zero.
The functional equation carries the result to the other side
of the critical line. See Section~\ref{sec:analytic}.

For local divisibility and the slope comparison in the original proof,
see the accounts by Slijepčević and Zhang. The analytic reduction described
by Slijepčević uses bounds on intervals \([Y,Y^{3/2}]\) with arbitrarily
large \(Y\), as in the original proof.
\cite[Section~1.3]{slijepcevic} \cite[Sections~1 and~4]{zhang}
